% =====================================================================
%  25.5  The Engine Boundary -- when a link-cut tree is the WRONG tool.
%  Exports: sec:c23-boundary  def:c23-boundary-pick
%  WORKED examples only -- NO solvedproblem (in-section solved stays 2).
%  Discharges the forward-cites promised by Ch.22 (sec:c16-boundary) and
%  Ch.24 (link-cut is the dynamic-topology neighbour). Mirrors the
%  ch18 05_boundary worked-only style.
% =====================================================================
\section{The Engine Boundary: When Link-Cut Is the Wrong Tool}\label{sec:c23-boundary}

\begin{previouslybox}[Every engine in this book ends by naming its own limits]
Chapter~22 closes its HLD chapter by listing the problems HLD must hand away
(\(\secref{c16-boundary}\)); Chapter~24 does the same for centroid decomposition.
An LCT is the most general of the three -- it can reroot, rewire and fold a moving
path -- and generality is exactly why it is the most expensive to write and the
easiest to get wrong. The mature reflex is not ``can an LCT solve this?'' (it can
solve almost anything about a forest) but ``should I use one?'' This section is a
worked reading of six specs and the cheaper tool each one actually wants.
\end{previouslybox}

\begin{topicintro}
A link-cut tree earns its constants only when the \emph{topology} changes between
queries. Fix the tree and every job it does has a cheaper specialist; keep only
connectivity and a DSU beats it outright; know all edits offline and a rollback DSU
or a divide-and-conquer over the timeline runs in a fraction of the code. The
section walks six short specifications through the pick
(\(\defref{c23-boundary-pick}\)); none has a solved problem, because the whole
lesson is \emph{not} writing an LCT.

\textbf{Prerequisites.} The whole of \(\secref{c23-splay}\)--\(\secref{c23-beyond}\),
so you know what an LCT is before you decide against it; \(\secref{c16-boundary}\)
(HLD's boundary), Ch.~5 (\chref{c24}) Euler tour, Ch.~17 (\chref{c28}) DSU, Ch.~24
(\chref{c18}) offline D\&C.

\textbf{After this section you can}
\begin{itemize}[nosep]
  \item read a dynamic-tree spec and spot whether the topology truly changes; \hfill\lv{5}
  \item route fixed-tree path/subtree jobs to HLD / Euler+Fenwick; \hfill\lv{5}
  \item route offline connectivity to rollback DSU and reserve the LCT for online path work. \hfill\lv{5}
\end{itemize}
\end{topicintro}

\begin{definition}{The dynamic-tree tool pick}{c23-boundary-pick}
Ask two questions in order. \textbf{(Q1) Does the topology change between
queries?} No \(\Rightarrow\) it is a \emph{static} tree: path range add/sum \(\to\)
HLD (\chref{c16}); subtree range \(\to\) Euler tour + Fenwick (\chref{c24}); a value
per node \(\to\) tree DP / rerooting (\chref{c10}, \chref{c20}). \textbf{(Q2) If it
does change, do you need a value \emph{along the path}, or only
\emph{connectivity}?} Path value, online \(\Rightarrow\) LCT
(\(\secref{c23-makeroot}\)); subtree value, online \(\Rightarrow\) LCT with virtual
children (\(\secref{c23-beyond}\)) or an Euler-tour tree; connectivity only, online
\(\Rightarrow\) Euler-tour tree; connectivity only, offline \(\Rightarrow\) rollback
DSU / segment-tree-on-time (\(\chref{c28}\), \(\chref{c18}\)).
\end{definition}

\begin{comparisonbox}[Same problem, three budgets: LCT vs HLD vs offline]
\setlength{\tabcolsep}{3pt}\renewcommand{\arraystretch}{1.2}%
\begin{tabular}{@{}p{0.29\linewidth}p{0.21\linewidth}p{0.21\linewidth}p{0.21\linewidth}@{}}
\toprule
\textbf{Job} & \textbf{Fixed tree} & \textbf{Changing, online} & \textbf{Changing, offline}\\
\midrule
path sum/min & HLD + segtree \(\bigO{\log^2 n}\) & LCT \(\bigO{\log n}\) amort & offline D\&C on time \\
subtree sum & Euler + Fenwick \(\bigO{\log n}\) & ETT or LCT + \(\mathit{vir}\) & rollback / D\&C \\
connectivity & trivial (one tree) & ETT or LCT \code{findroot} & rollback DSU \(\bigO{\alpha}\) \\
\bottomrule
\end{tabular}
\end{comparisonbox}

\subsection*{Six specs, read and routed}

\begin{itemize}[leftmargin=1.4em]
  \item \textbf{``Fixed tree, path add and path sum, \(2\cdot10^5\) queries.''} No
        edge ever changes. \(\Rightarrow\) \emph{HLD} (\chref{c16}) -- an LCT would
        pass but with a larger constant and triple the code. \textbf{Do not.}
  \item \textbf{``Fixed tree, subtree update and subtree sum.''} \(\Rightarrow\)
        \emph{Euler tour + Fenwick} (\chref{c24}) in one line of insight. \textbf{Do not.}
  \item \textbf{``Add/remove edges; only ask if \(u,v\) connected; all queries known
        in advance.''} \(\Rightarrow\) \emph{rollback DSU} or D\&C on the edge-activity
        intervals (\chref{c18}). \textbf{Do not.}
  \item \textbf{``Link/cut online; only connectivity.''} \(\Rightarrow\) \emph{Euler-tour
        tree} (or an LCT \code{findroot}); the ETT is less error-prone when no path
        value is needed. \textbf{Prefer ETT.}
  \item \textbf{``Link/cut online; sum along the current path.''} \(\Rightarrow\) \textbf{LCT.}
        The genuine, textbook use (\(\secref{c23-makeroot}\)).
  \item \textbf{``Over all simple paths in a changing \emph{graph}, the bottleneck
        value.''} \(\Rightarrow\) \textbf{LCT}, but only after a condensation (block-cut
        or spanning-forest) turns it into a forest -- which is precisely the Summit.
\end{itemize}

\begin{insightbox}[The honest cost of generality]
An LCT's \code{access}/\code{splay} hides a big constant and, worse, a large
\emph{bug} surface: lazy reverse tags, virtual-child bookkeeping, \code{pushdown}
order. A fixed-tree HLD gives you the same asymptotics on \emph{no} topology change
with array-based segment trees you can debug in minutes. The reason this chapter
exists is the fifth and sixth specs -- the ones HLD genuinely cannot serve. Treat a
topology that never changes as the loudest possible signal that you wanted a static
engine.
\end{insightbox}

\begin{warningbox}[The four ways students reach for an LCT too early]
\begin{enumerate}[nosep,leftmargin=1.6em]
  \item ``the problem mentions a dynamic tree'' but the edges are fixed and only
        \emph{vertex values} change -- that is HLD/Euler, not LCT.
  \item connectivity-only, asked offline -- a rollback DSU is \(\bigO{\alpha}\) with a
        tenth of the code (\chref{c28}).
  \item counting over \emph{all} paths of a fixed tree -- centroid decomposition
        (\chref{c18}), not a per-path fold.
  \item a state per node with no route dimension -- tree DP / rerooting (\chref{c10}).
\end{enumerate}
Each wastes an LCT on a job a simpler engine does faster and correctly.
\end{warningbox}

\begin{hardeningbox}
\begin{itemize}[nosep,leftmargin=1.4em]
  \item \textbf{Prove the topology moves.} Before coding, find the exact operation that
        adds or deletes an \emph{edge}. If you cannot, you have an HLD/Fenwick problem.
  \item \textbf{Check the constant.} With \(n,q\le 2\cdot10^5\) and a fixed tree, an
        LCT's \(\sim\)20 splay-steps per query versus HLD's segment-tree walk can be the
        difference between AC and TLE even though both are ``\(O(\log n)\).''
  \item \textbf{Offline \(\Rightarrow\) no online structure.} If you may read all queries
        first, prefer rollback DSU / D\&C; the LCT is the tool for when you \emph{cannot}.
\end{itemize}
\end{hardeningbox}

\begin{recallbox}[25.5 in one breath]
\emph{An LCT is for a \textbf{changing topology} with a \textbf{path (or subtree) value}
query, \textbf{online}. Fixed tree \(\to\) HLD / Euler+Fenwick; connectivity only \(\to\)
ETT or rollback DSU; all-pairs counting \(\to\) centroid; per-node state \(\to\) tree DP.
Two questions decide it: does an edge ever appear or vanish, and do you need a value
along the route? Both ``yes'' \(\Rightarrow\) LCT; otherwise hand it off
(\(\defref{c23-boundary-pick}\)).}
\selftest{
  \item A spec has path sums but the tree never changes -- which engine? \emph{(HLD, \chref{c16}; the LCT's generality is dead weight.)}
  \item Connectivity is all you need and queries are offline -- best tool? \emph{(Rollback DSU / D\&C on time, \chref{c28}/\chref{c18}.)}
  \item State the two questions of the pick. \emph{(Does the topology change between queries? Do you need a value along the path? \(\defref{c23-boundary-pick}\).)}
}
\end{recallbox}
